\documentclass[11pt,a4paper]{article}
% Shared preamble for the English probability and statistics documents.
% Based on the original course preamble. Compile with pdfLaTeX.
\usepackage[top=1.5cm,bottom=2cm,left=2cm,right=2cm]{geometry}
\usepackage[T1]{fontenc}
\usepackage{lmodern}
\usepackage[utf8]{inputenc}
\usepackage[english]{babel}

\usepackage{amsmath}
\usepackage{amssymb}
\usepackage{amsthm}
\usepackage{mathrsfs}
\usepackage{mathtools}
\usepackage{etoolbox}
\usepackage{tikz}
\usetikzlibrary{shapes.multipart,arrows.meta}
\usepackage{xcolor}
\usepackage{enumerate}
\usepackage[shortlabels]{enumitem}
\usepackage{hyperref}
\usepackage{version}
\usepackage{bbold}
\usepackage{fontawesome5}
\usepackage{fancyhdr}
\lfoot{}
\usepackage{natbib}
\usepackage{graphicx}
\usepackage{pdfpages}
\usepackage{authblk}
    
\usepackage{blindtext}
\usepackage{multicol}
\usepackage{collcell}
\usepackage{wrapfig}
\usepackage{blkarray}
\usepackage{cancel}
\usepackage{appendix}
\usepackage{float}
\makeatletter
\input{tex/assets/slashbox.sty}
\makeatother
\hypersetup{unicode,colorlinks,citecolor=red,filecolor=black,
  linkcolor=red,urlcolor=cyan}
\setlength{\columnsep}{1cm}
\pagenumbering{gobble}

\def\numtext#1{%
  \ifcase#1\or one\or two\or three\or four\or five\or six\or seven\or eight\or nine\or ten\or
  eleven\or twelve\or thirteen\or fourteen\or fifteen\or sixteen\or seventeen\or eighteen\or nineteen\or twenty\or Lots\fi}

\newcommand{\Byte}[1]{%
  \def\nodecontents{}%
  \foreach[count=\l] \k in {#1} {\xappto{\nodecontents}{\noexpand\nodepart{\numtext{\l}}\k}}%
  \node[rectangle split,rectangle split parts=8,rectangle split horizontal,draw] at (2,2){\nodecontents};}

\newcommand{\D}{{\mathrm{d}}}
\newcommand{\cE}{{\mathcal{E}}}
\newcommand{\cG}{{\mathcal{G}}}
\newcommand{\cX}{{\mathcal{X}}}
\newcommand{\cY}{{\mathcal{Y}}}
\newcommand{\cB}{{\mathcal{B}}}
\newcommand{\cF}{{\mathcal{F}}}
\newcommand{\cO}{{\mathcal{O}}}
\newcommand{\cS}{{\mathcal{S}}}
\newcommand{\cT}{{\mathcal{T}}}
\newcommand{\given}{\, \big| \,}
\newcommand{\cL}{{\mathscr{L}}}
\newcommand{\cD}{{\mathcal{D}}}
\newcommand{\cA}{{\mathscr{A}}}
\newcommand{\cP}{{\mathscr{P}}}
\newcommand{\cR}{{\mathscr{R}}}
\newcommand{\A}{{\mathcal A}}
\newcommand{\one}{\mathbb{1}}
\newcommand{\mathbbm}[1]{\mathbb{#1}}

\providecommand{\C}{}
\renewcommand{\C}{\mathbb C}
\newcommand{\R}{\mathbb R}
\newcommand{\Q}{\mathbb Q}
\newcommand{\N}{\mathbb N}
\newcommand{\T}{\mathbb T}
\newcommand{\F}{\mathbb F}
\renewcommand{\P}{\mathbb P}
\newcommand{\Z}{\mathbb Z}
\newcommand{\E}{\mathbb {E}}
\newcommand{\Pbb}[1]{\P\left( #1\right)}
\newcommand{\Ebb}[1]{\E\left[ #1\right]}

\newcommand{\lb}{\left (}
\newcommand{\rb}{\right )}
\newcommand{\lbb}{\left [}
\newcommand{\rbb}{\right ]}
\newcommand{\labs}{\left |}
\newcommand{\rabs}{\right |}
\newcommand{\lbrb}[1]{\lb #1 \rb}
\newcommand{\lbbrbb}[1]{\lbb#1\rbb}
\newcommand{\labsrabs}[1]{\labs#1\rabs}
\newcommand{\lbbrb}[1]{\lbb#1\rb}
\newcommand{\lbrbb}[1]{\lb#1\rbb}
\newcommand{\langlerangle}[1]{\langle#1\rangle}
\newcommand{\lbcurly}{\{}
\newcommand{\rbcurly}{\}}
\newcommand{\curly}[1]{\lbcurly#1\rbcurly}
\newcommand{\abs}{\labsrabs}
\newcommand{\ind}[1]{1_{#1}}
\newcommand{\indep}{\perp \!\!\! \perp}
\renewcommand{\d}{\, \mathrm{d}}

\newtheorem{theorem}{Theorem}
\newtheorem{problemCnt}{Lemma}
\theoremstyle{definition}
\newtheorem{problem}[problemCnt]{Problem}
\newtheorem{problembook}[problemCnt]{\faBook \,  Problem}
\newtheorem{problempencil}[problemCnt]{\faPencil* \,  Problem}
\newtheorem{problemstar}[problemCnt]{\faStar\,  Problem}
\newtheorem{problemstarpencil}[problemCnt]{\faStar \, \faPencil* \,   Problem}
\newtheorem{probleminteresting}[problemCnt]{\faSmile[regular] \,   Problem}
\newtheorem{solution}{Solution}

\newcommand{\fwcell}[1]{\makebox[\arraystretch\normalbaselineskip]{$#1$}}
\newenvironment{enumeratei}{\begin{enumerate}[(i)]}{\end{enumerate}}
\newenvironment{enumeratea}{\begin{enumerate}[(a)]}{\end{enumerate}}
\newenvironment{enumeraten}{\begin{enumerate}[\upshape 1.]}{\end{enumerate}}

% Compatibility helpers needed by a few files in the combined archive.
\AtBeginDocument{\noindent{\footnotesize English translation prepared with ChatGPT (OpenAI, GPT-6), 26 October 2026, for translation of the Bulgarian text into English and LaTeX formatting.
The Bulgarian originals remain the reference; translation may contain errors.}\par\smallskip}
\newcommand{\MissingGraphic}[1]{%
  \begin{center}\fbox{\parbox{0.82\linewidth}{\centering Missing source image: \texttt{#1}}}\end{center}}



\pagestyle{fancy}\fancyhf{}\lhead{Computer Science: Test 2 solutions}\rhead{8 June 2024}
\setlength{\headheight}{15pt}\pagenumbering{arabic}\cfoot{\thepage}
\begin{document}
\noindent\emph{English transcription of the handwritten solution notes. Original calculations are retained; apparent inconsistencies are identified in the transcription notes.}

\section*{Problem 1}
We have $\E X_i=5011$ and $D(X_i)=4000$.
\subsection*{1. The sum}
Let $S_n=\sum_{i=1}^nX_i$. We seek $n$ such that $\P(S_n>10^6)\ge99\%$. By the central limit theorem,
\begin{align*}
\P(S_n>10^6)
&=\P\left(\frac{S_n-\E S_n}{\sqrt{D(S_n)}}>
\frac{10^6-5011n}{\sqrt{4000n}}\right)\\
&\approx\P\left(N(0,1)>\frac{10^6-5011n}{\sqrt{4000n}}\right).
\end{align*}
For this to be at least $99\%$, require
\[
\frac{10^6-5011n}{\sqrt{4000n}}<-2.33,
\qquad\text{giving}\qquad n\ge200.
\]

\subsection*{2. The maximum}
We seek $n$ such that
\[
\P\left(\max_{1\le i\le n}X_i>10^4\right)\ge99\%.
\]
Equivalently,
\[
\P\left(\max_{1\le i\le n}X_i\le10^4\right)\le1\%,
\qquad [\P(X_1\le10^4)]^n\le1\%.
\]
Writing $p=\P(X_1\le10^4)$, this gives
\[
n\ge\frac{\ln(1/100)}{\ln p}.
\]
To obtain a particular value of $p$, a specific distribution of $X_i$ must be chosen. If $X_i\sim N(5011,4000)$, then
\[
\P(X_1\le10^4)=\P\left(N(0,1)\le\frac{10^4-5011}{\sqrt{4000}}\right)
\approx\P(N(0,1)\le78.88).
\]
Thus the probability for the maximum involves $\Phi(78.88)^n$.
All solutions that reach an expression of this kind receive the maximum possible number of points.

\subsection*{Bonus: an explicit bound}
To find an explicit value of $n$ such that $\Phi(78.88)^n\le1\%$, note that although $\Phi(78.88)$ is very close to 1, raising it to a sufficiently high power makes it close to 0.
The standard bounds for $\P(N(0,1)>x)=1-\Phi(x)$ are
\[
\left(\frac1x-\frac1{x^3}\right)\frac{e^{-x^2/2}}{\sqrt{2\pi}}
\le1-\Phi(x)\le\frac1x\frac{e^{-x^2/2}}{\sqrt{2\pi}}.
\]
\paragraph{Proof.} Integrate by parts twice:
\begin{align*}
1-\Phi(x)&=\int_x^\infty\frac1{\sqrt{2\pi}}e^{-t^2/2}\d t\\
&=\frac1{\sqrt{2\pi}}\left(-\int_x^\infty\frac1t\d(e^{-t^2/2})\right)\\
&=\frac1{\sqrt{2\pi}}\left(\frac1x e^{-x^2/2}
+\int_x^\infty e^{-t^2/2}\d(1/t)\right)\\
&=\frac1{\sqrt{2\pi}}\left(\frac1x e^{-x^2/2}
-\int_x^\infty\frac1{t^2}e^{-t^2/2}\d t\right)\\
&=\frac1{\sqrt{2\pi}}\left(\frac1x e^{-x^2/2}
+\int_x^\infty\frac1{t^3}\d(e^{-t^2/2})\right)\\
&=\frac1{\sqrt{2\pi}}\left(\frac1x e^{-x^2/2}-\frac1{x^3}e^{-x^2/2}
-\int_x^\infty e^{-t^2/2}\d(1/t^3)\right).
\end{align*}
The final line in the original is written
\[
=\frac1{\sqrt{2\pi}}\left(\frac1x-\frac1{x^3}\right)e^{-x^2/2}
-\frac3{\sqrt{2\pi}}\int_x^\infty\frac1{t^4}e^{-t^2/2}\d t
\ge\frac1{\sqrt{2\pi}}\left(\frac1x-\frac1{x^3}\right)e^{-x^2/2}.\tag{Note 1}
\]
Hence, for $x=78.88$,
\[
\Phi(78.88)\le1-\left(\frac1x-\frac1{x^3}\right)
\frac{e^{-x^2/2}}{\sqrt{2\pi}}.
\]
It is sufficient to find $n$ such that
\[
\left[1-\left(\frac1x-\frac1{x^3}\right)
\frac{e^{-x^2/2}}{\sqrt{2\pi}}\right]^n\le1\%.
\]
The source next writes
\[
n\ln\left[1-\left(\frac1x-\frac1{x^3}\right)
\frac{e^{-x^2/2}}{\sqrt{2\pi}}\right]\le1\%.\tag{Note 2}
\]
Using $\ln(1+u)\le u$ (or $\ln(1-u)\sim-u$), the subsequent line is
\[
n\left[-\left(\frac1x-\frac1{x^3}\right)
\frac{e^{-x^2/2}}{\sqrt{2\pi}}\right]\le\ln(1/100).
\]
This gives the sufficient bound
\[
n\ge\ln(100)e^{x^2/2}\sqrt{2\pi}\,\frac{x^3}{x^2-1},
\qquad\text{and thus}\qquad n\ge1.16\cdot10^{1354}!
\]

\section*{Problem 2}
See Problem 3 of the Software Engineering examination, January 2020/2021.

\section*{Problem 3}
Let $X\sim\operatorname{Exp}(\alpha^{-1})$ and $Y\sim\operatorname{Exp}(\beta^{-1})$.
\subsection*{1. Estimating the mean}
We seek $c$ such that $\P(|\alpha-\overline X|>c)\le1\%$.
The handwritten source gives the following CLT normalization for large $n$:
\[
\frac{\overline X-\alpha}{\alpha\sqrt n}\approx Z\sim N(0,1).\tag{Note 3}
\]
It therefore seeks $c$ such that
\[
\P\left(|Z|>\frac c{\alpha\sqrt n}\right)\le1\%
\quad\Longleftrightarrow\quad
\P\left(Z>\frac c{\alpha\sqrt n}\right)\le0.5\%.
\]
The source concludes
\[
\frac c{\alpha\sqrt n}\approx2.57,
\qquad c\approx2.57\alpha\sqrt n.\tag{Note 3}
\]

\subsection*{2. The ratio}
Set $Z_i=\sqrt{X_i/Y_i}$. If $\E Z_i<\infty$, the strong law of large numbers gives
\[
\frac{Z_1+\cdots+Z_n}{n}\longrightarrow\E Z_1
\qquad\text{almost surely}.
\]
\paragraph{First method.} For $t>0$, using $\P(\operatorname{Exp}(\lambda)>t)=e^{-\lambda t}$,
\begin{align*}
\P\left(\sqrt{X/Y}>t\right)
&=\P(X>Yt^2)=\int_0^\infty\P(X>yt^2)f_Y(y)\d y\\
&=\int_0^\infty e^{-yt^2/\alpha}\frac1\beta e^{-y/\beta}\d y\\
&=\left(\frac{t^2}\alpha+\frac1\beta\right)^{-1}\frac1\beta
\int_0^\infty e^{-y(t^2/\alpha+1/\beta)}\d\left[y(t^2/\alpha+1/\beta)\right]\\
&=\frac{\alpha\beta}{\beta t^2+\alpha}\frac1\beta
=\frac\alpha{\beta t^2+\alpha}.
\end{align*}
The integral after substitution equals $\int_0^\infty e^{-x}\d x=1$. Therefore,
\begin{align*}
\E\sqrt{X/Y}
&=\int_0^\infty\P\left(\sqrt{X/Y}>t\right)\d t\\
&=\int_0^\infty\frac\alpha{\beta t^2+\alpha}\d t
=\int_0^\infty\frac1{(\beta/\alpha)t^2+1}\d t\\
&=\sqrt{\frac\alpha\beta}\int_0^\infty\frac{\d y}{1+y^2}
=\sqrt{\frac\alpha\beta}\frac\pi2,
\end{align*}
where $y=\sqrt{\beta/\alpha}\,t$. This is different from
$\sqrt{\E X/\E Y}=\sqrt{\alpha/\beta}$.

\paragraph{Second method.} Independence gives
\[
f_{X,Y}(x,y)=f_X(x)f_Y(y)=\frac1\alpha e^{-x/\alpha}\frac1\beta e^{-y/\beta}.
\]
Thus
\begin{align*}
\E\sqrt{X/Y}
&=\int_0^\infty\int_0^\infty\sqrt{x/y}\,
\frac1\alpha e^{-x/\alpha}\frac1\beta e^{-y/\beta}\d x\d y\\
&=\left(\int_0^\infty\sqrt x\frac1\alpha e^{-x/\alpha}\d x\right)
\left(\int_0^\infty\frac1{\sqrt y}\frac1\beta e^{-y/\beta}\d y\right).
\end{align*}
Substitute $x'=x/\alpha$ and $y'=y/\beta$. The factors $\sqrt\alpha$ and $1/\sqrt\beta$ can be taken outside, giving
\[
\E\sqrt{X/Y}=\sqrt{\frac\alpha\beta}\Gamma(3/2)\Gamma(1/2)
=\sqrt{\frac\alpha\beta}\frac12\Gamma(1/2)^2
=\sqrt{\frac\alpha\beta}\frac\pi2.
\]
Here $\Gamma(z)=\int_0^\infty x^{z-1}e^{-x}\d x$ and $\Gamma(z+1)=z\Gamma(z)$.

\paragraph{Third method.} Introduce $Z=\sqrt{X/Y}$ and $W=\sqrt Y$, so $X=Z^2W^2$ and $Y=W^2$. The Jacobian is
\[
|J|=\begin{vmatrix}2zw^2&2z^2w\\0&2w\end{vmatrix}=4zw^3.
\]
Thus, for $z,w>0$,
\[
f_{Z,W}(z,w)=f_{X,Y}(z^2w^2,w^2)\,4zw^3
=\frac1{\alpha\beta}e^{-z^2w^2/\alpha-w^2/\beta}\,4zw^3.
\]
Integrating out $w$,
\begin{align*}
f_Z(z)&=\int_0^\infty\frac1{\alpha\beta}
e^{-w^2(1/\beta+z^2/\alpha)}\,4zw^3\d w\\
&=\frac{4z}{\alpha\beta}\frac12\int_0^\infty
e^{-w^2(1/\beta+z^2/\alpha)}w^2\d(w^2)\\
&=\frac{2z}{\alpha\beta}\left(\frac1\beta+\frac{z^2}\alpha\right)^{-2}
\int_0^\infty e^{-y}y\d y\\
&=\frac{2z}{\alpha\beta}\frac{(\alpha\beta)^2}{(\alpha+\beta z^2)^2}
=\frac{2\alpha\beta z}{(\alpha+\beta z^2)^2}
=2\frac\beta\alpha\frac z{(1+(\beta/\alpha)z^2)^2}.
\end{align*}
We used $y=w^2(1/\beta+z^2/\alpha)$ and $\int_0^\infty ye^{-y}\d y=\E[\operatorname{Exp}(1)]=1$.
Finally, with $z=\sqrt{\alpha/\beta}\,y$,
\begin{align*}
\E Z&=\int_0^\infty z\,2\frac\beta\alpha\frac z{(1+(\beta/\alpha)z^2)^2}\d z\\
&=2\sqrt{\frac\alpha\beta}\int_0^\infty\frac{y^2}{(1+y^2)^2}\d y
=2\sqrt{\frac\alpha\beta}\frac\pi4
=\frac\pi2\sqrt{\frac\alpha\beta}.
\end{align*}

\section*{Problem 4}
See Problem 4 of the Software Engineering examination, February 2024.

\section*{Transcription notes}
\begin{enumerate}
\item The remainder after the second integration by parts should have a plus sign. The minus sign in the original conflicts with its next inequality.
\item After taking logarithms, the right-hand side should be $\ln(1/100)$ rather than $1\%$. The original's following line uses the logarithm.
\item The original repeatedly writes $\alpha\sqrt n$ for the standard deviation of the sample mean. For $X_i\sim\operatorname{Exp}(1/\alpha)$ it is $\alpha/\sqrt n$. The corresponding approximate bound is $c\approx2.57\alpha/\sqrt n$.
\end{enumerate}
\end{document}
